\section{Comparable prime supply}
\label{app:prime-supply}

This appendix proves the ordinary analytic input used in Sections~\ref{sec:signed-inverse} and~\ref{sec:scalar-finality}: one fixed-ratio prime band containing $\asymp X/\log X$ primes.  The multiplier is chosen to be an integer so that every later cutoff remains on the natural-number scale.

Put
\[
\theta(x)=\sum_{p\le x}\log p,
\qquad
\psi(x)=\sum_{p^m\le x}\log p
=\log\operatorname{lcm}(1,\ldots,\lfloor x\rfloor).
\]
The central binomial coefficient satisfies
\[
\binom{2n}{n}\mid\operatorname{lcm}(1,\ldots,2n+1).
\]
Indeed, for each prime $p$, Legendre's formula gives
\[
v_p\binom{2n}{n}
=\sum_{j\ge1}
\left(\left\lfloor\frac{2n}{p^j}\right\rfloor
-2\left\lfloor\frac n{p^j}\right\rfloor\right),
\]
and every summand is $0$ or $1$; if the sum is $e$, then some nonzero summand has index at least $e$, so $p^e\le2n$.

Since the central binomial coefficient is maximal in its row,
\[
4^n\le(2n+1)\binom{2n}{n}.
\]
The divisibility above therefore gives
\[
\psi(2n+1)
\ge\log\binom{2n}{n}
\ge n\log4-\log(2n+1).
\]
Taking $n=\lfloor(x-1)/2\rfloor$ and using monotonicity of $\psi$ shows that, for some $c_\psi>0$,
\[
\psi(x)\ge c_\psi x
\]
for all sufficiently large $x$.

On the other hand, every prime $n<p\le2n$ divides $\binom{2n}{n}$, so
\[
\theta(2n)-\theta(n)
\le\log\binom{2n}{n}
\le2n\log2.
\]
Dyadic telescoping gives $\theta(x)\le C_\theta x$ for a fixed $C_\theta$.  Since
\[
\psi(x)=\sum_{m\ge1}\theta(x^{1/m}),
\]
there are only $O(\log x)$ nonzero terms with $m\ge2$, and the upper bound for $\theta$ gives
\[
0\le\psi(x)-\theta(x)
=\sum_{m\ge2}\theta(x^{1/m})
=O(\sqrt x\log x)
=o(x).
\]
Subtracting this from the lower bound for $\psi$ yields, after increasing the threshold if necessary,
\[
c_\theta x\le\theta(x)\le C_\theta x
\tag{P1}\label{eq:chebyshev-linear}
\]
for some $c_\theta>0$.

Choose an integer $\Lambda\ge2$ such that $C_\theta<c_\theta\Lambda$.  Then
\[
\theta(\Lambda X)-\theta(X)
\ge(c_\theta\Lambda-C_\theta)X,
\]
while every prime in $(X,\Lambda X]$ contributes between $\log X$ and $\log(\Lambda X)$.  Thus:

\begin{proposition}[Comparable prime band]\label{prop:prime-band}
There are an integer $\Lambda\ge2$ and constants $c,C>0$ such that, for every sufficiently large integer $X$,
\[
c\frac X{\log X}
\le
\#\{p:X<p\le\Lambda X\}
\le
C\frac X{\log X}.
\]
\end{proposition}

Only the existence of fixed constants $\Lambda,c,C$ is used in the main argument; their numerical values play no role.
